% ============================================================================
% COURS 03 — EXERCICES DE TRANSMETTEURS
% Source : 4-ModelisationTransmetteurs2022.docx
% ============================================================================

\section{Exercices — modélisation des transmetteurs}
\label{sec:ex-trans}

\subsection{Réducteurs épicycloïdaux ATV et Pellenc}

\textit{Niveau 2.}

Deux réducteurs compacts sont étudiés à partir de leurs schémas
cinématiques.

\begin{center}
\TLFourImage[.45\linewidth]{image81.png}\hfill
\TLFourImage[.45\linewidth]{image82.png}
\end{center}

Pour le réducteur ATV :
\[
Z_1=166,\quad Z_2=160,\quad Z_{2'}=164,\quad Z_3=170.
\]
Pour le sécateur Pellenc :
\[
Z_1=19,\quad Z_2=19,\quad Z_3=57.
\]

\begin{enumerate}
  \item Identifier les planétaires, le porte-satellite et les satellites.
  \item Déterminer la raison de base $\lambda$ de chaque train.
  \item Écrire la formule de Willis.
  \item Introduire l'élément bloqué et calculer le rapport de réduction.
  \item Vérifier que le résultat correspond à un fort rapport dans un encombrement réduit.
\end{enumerate}

\subsection{Train d'engrenages simple}

\textit{Niveau 1.}

Un train d'engrenages à axes parallèles comporte :
\[
Z_1=65,\ Z_2=32,\ Z_3=24,\ Z_{3'}=48,\ Z_4=38,
\]
\[
Z_{4'}=82,\ Z_5=26,\ Z_{5'}=54,\ Z_6=42,\ Z_7=30.
\]

\TLFourImage[.72\linewidth]{image83.png}

\begin{enumerate}
  \item Indiquer le sens de rotation de chaque roue.
  \item Lister les roues menantes et les roues menées.
  \item Donner l'expression du rapport global.
  \item Faire l'application numérique et conclure : réducteur ou multiplicateur.
  \item Identifier les roues folles et expliquer leur effet sur le rapport et le signe.
\end{enumerate}

\subsection{Monte-charge DEMAG}

\textit{Niveau 2.}

Un moteur tournant à $1500\ \mathrm{tr\,min^{-1}}$ entraîne un tambour
par un réducteur. Les roues sont caractérisées par :

\begin{center}
\small
\begin{tabular}{c|cccccc}
Repère & 6 & 9a & 9b & 11a & 11b & 16\\ \hline
Module & 1 & 1 & 1 & 1 & 1,25 & 1,25\\
Dents & 16 & 46 & 19 & 59 & 17 & 85
\end{tabular}
\end{center}

\begin{center}
\TLFourImage[.42\linewidth]{image85.png}\hfill
\TLFourImage[.50\linewidth]{image90.jpeg}
\end{center}

\begin{enumerate}
  \item Identifier et colorier les classes d'équivalence cinématique.
  \item Dessiner le schéma cinématique dans le plan de la coupe C--C.
  \item Établir le rapport de réduction du train simple puis le calculer.
  \item Déterminer la vitesse du tambour et la vitesse de levage.
  \item Une vis sans fin à quatre filets doit fournir le même rapport. Déterminer le nombre de dents de la roue associée.
\end{enumerate}

\subsection{Portail automatisé}

\textit{Niveau 2.}

Le moteur 21 tourne à $3000\ \mathrm{tr\,min^{-1}}$ et entraîne le bras
de commande par quatre engrènements : 19--5, 3--22, 23--17 et 16--11.
Les modules sont respectivement $0{,}75$, $0{,}75$, $1$ et $1{,}5$.

\begin{center}
\TLFourImage[.31\linewidth]{image91.png}\hfill
\TLFourImage[.31\linewidth]{image93.png}\hfill
\TLFourImage[.31\linewidth]{image94.jpeg}
\end{center}

\begin{enumerate}
  \item Déterminer le rapport global en fonction des diamètres primitifs.
  \item Remplacer les diamètres par les nombres de dents lorsque les modules sont identiques dans un engrènement.
  \item Calculer la vitesse de l'arbre 9.
  \item En déduire le temps nécessaire pour parcourir $90^\circ$.
  \item Conclure sur le respect du critère de temps d'ouverture.
\end{enumerate}

\subsection{Roue d'engin à deux trains épicycloïdaux}

\textit{Niveau 4.}

La pièce 7 est liée au châssis, l'arbre 1 est entraîné par un moteur
hydraulique et la roue est fixée sur la pièce 12.

\begin{center}
\TLFourImage[.43\linewidth]{image97.jpeg}\hfill
\TLFourImage[.43\linewidth]{image98.jpeg}
\end{center}

Données :
\begin{itemize}
  \item train 1 : $Z_1=21$, $Z_2=51$, $Z_{12a}=123$ ;
  \item train 2 : $Z_4=23$, $Z_5=34$, $Z_{12b}=91$.
\end{itemize}

\begin{enumerate}
  \item Établir la relation de Willis du premier train entre $\omega_{1/7}$, $\omega_{4/7}$ et $\omega_{12/7}$.
  \item Établir la relation du second train entre $\omega_{4/7}$ et $\omega_{12/7}$.
  \item Éliminer $\omega_{4/7}$ et déterminer $\omega_{12/7}/\omega_{1/7}$.
  \item Vérifier le sens de rotation de la roue et l'ordre de grandeur de la réduction.
\end{enumerate}

\subsection{Poulie REDEX}

\textit{Niveau 3.}

\TLFourImage[.66\linewidth]{image99.png}

\begin{enumerate}
  \item Repérer les différents étages du réducteur.
  \item Identifier les roues solidaires et les contacts intérieurs ou extérieurs.
  \item Déterminer le rapport $\omega_{31/0}/\omega_{5/0}$.
  \item Vérifier la cohérence du signe et du caractère réducteur.
\end{enumerate}

\subsection{Robot ERICC}

\textit{Niveau 2.}

La poulie motrice 18 est liée à la sortie du moto-réducteur. La poulie
23 est fixe par rapport au bâti ; le socle chaise 10 tourne donc autour
de l'axe de lacet.

\begin{center}
\TLFourImage[.31\linewidth]{image100.png}\hfill
\TLFourImage[.31\linewidth]{image101.jpeg}\hfill
\TLFourImage[.31\linewidth]{image102.png}
\end{center}

\begin{enumerate}
  \item Écrire la condition de roulement sans glissement de la courroie par rapport au solide 10.
  \item Déterminer $\omega_{23/10}/\omega_{18/10}$ en fonction des nombres de dents.
  \item Utiliser la composition des vitesses avec $\omega_{23/0}=0$.
  \item En déduire la vitesse de lacet du robot lorsque le moto-réducteur tourne à $50\ \mathrm{tr\,min^{-1}}$.
\end{enumerate}

\subsection{Dessileuse-pailleuse-distributrice}

\textit{Niveau 2.}

La boîte de vitesses réduit la fréquence de rotation afin d'augmenter le
couple transmis à la turbine de distribution d'ensilage.

\begin{center}
\TLFourImage[.46\linewidth]{image103.jpeg}\hfill
\TLFourImage[.46\linewidth]{image104.jpeg}
\end{center}

\begin{enumerate}
  \item Expliquer la solution constructive de sélection de la petite vitesse.
  \item Construire le schéma cinématique minimal dans cette configuration.
  \item Décrire le guidage en rotation de l'arbre 7 et les arrêts axiaux.
  \item Calculer le rapport $\omega_s/\omega_e$.
  \item Qualifier la transmission et discuter son rendement.
\end{enumerate}

\subsection{Coffre motorisé de voiture}

\textit{Niveau 2.}

Un moteur de $150\ \mathrm W$ tournant à $3300\ \mathrm{tr\,min^{-1}}$
entraîne le hayon par un réducteur puis un mécanisme à quatre barres. La
manivelle 4 doit parcourir $68{,}4^\circ$ ; le temps d'ouverture imposé
est compris entre $3$ et $5\ \mathrm s$.

\begin{center}
\TLFourImage[.30\linewidth]{image105.png}\hfill
\TLFourImage[.30\linewidth]{image106.png}\hfill
\TLFourImage[.34\linewidth]{image108.png}
\end{center}

\begin{enumerate}
  \item Identifier l'étage roue-vis sans fin et les engrenages du réducteur.
  \item Déterminer $i=\omega_{4/0}/\omega_{2/0}$.
  \item Calculer la vitesse de la manivelle 4 puis le temps d'ouverture.
  \item Conclure sur le cahier des charges.
  \item Discuter l'intérêt de l'irréversibilité lorsque le moteur n'est plus alimenté.
\end{enumerate}

\subsection{Objectif photographique et résolution des codeurs}

L'objectif comporte un codeur absolu linéaire 4 bits et un codeur
incrémental monovoie délivrant 30 impulsions par tour de moteur.

\begin{center}
\TLFourImage[.29\linewidth]{image8.png}\hfill
\TLFourImage[.29\linewidth]{image55.png}\hfill
\TLFourImage[.36\linewidth]{image56.png}
\end{center}

\begin{enumerate}
  \item Déterminer la relation entre le déplacement $d_l$ de la lentille et l'angle $\theta_{rc}$ du codeur incrémental.
  \item En déduire le déplacement correspondant à une impulsion sur un front montant.
  \item Comparer la résolution obtenue à celle du codeur absolu 4 bits sur la course totale.
  \item Identifier les avantages respectifs des deux technologies pour la prise d'origine et la précision.
\end{enumerate}
